% s06-sec3c.tex —— §3.5 沿稳定共轭类的积分（p.34–38）
\subsection{沿稳定共轭类的积分}
\srcpage{34}

存在 $A$ 的一个稠密开子集 $A^{\mathrm{rs}}$，使得对任何 $g\in G$，只要其像 $a\in A^{\mathrm{rs}}$，$g$ 的中心化子就是环面。对任何元素 $a\in A^{\mathrm{rs}}(F_v)$，$G$ 在 $a$ 上方的 $F_v$-点集（由 Kottwitz [K1]，它非空）构成一个稳定共轭类。集合 $A^{\mathrm{rs}}(F_v)$ 就是强正则半单稳定共轭类的集合。

本节的目标是把积分
\[
\int_{G(F_v)}f_v(g_v)\,dg_v
\]
表示为先沿 $c:G\to A$ 的纤维积分、再在 $A(F_v)$ 上积分。本文最后一部分需要由此得到的公式。首先要以相容的方式固定测度。我们按 3.2 节的定义选取测度。特别地，固定 $G$ 上定义在 $F$ 上的体积形式 $\omega_G$ 与 $A$ 上定义在 $F$ 上的体积形式 $\omega_A$。配上加法特征的选取，就得到 $G(F_v)$ 上的测度 $|\omega_G|_v$ 与 $A(F_v)$ 上的测度 $|\omega_A|_v$。我们还有归一化测度 $dg_v=L_v(1,\sigma_G)|\omega_G|_v$ 与 $db_v=L_v(1,\sigma_Z)|\omega_A|_v$，归一化因子 $L_v(1,\sigma_G)=L_v(1,\sigma_Z)$ 相同。

第一个观察是，可以只限于 $A^{\mathrm{rs}}(F_v)$ 与 $G^{\mathrm{rs}}(F_v)$，其中 $G^{\mathrm{rs}}$ 是 $A^{\mathrm{rs}}$ 的原像。

\setcounter{theorem}{25}
\begin{lemma}
设 $f_v$ 是 $G(F_v)$ 上光滑紧支函数，$dg$ 是 $G(F_v)$ 上的不变测度。用同样的记号表示它们在 $G^{\mathrm{rs}}(F_v)$ 上的限制。则积分
\[
\int_{G^{\mathrm{rs}}(F_v)}f_v(g_v)\,dg_v
\]
绝对收敛，且等于 $\int_{G(F_v)}f_v(g_v)\,dg_v$。
\end{lemma}

这来自一个一般事实：严格闭子概形的 $F_v$-点集测度为零。\footnote{我们本想给出一个或几个参考文献。这个引理的证明应当远比关于"严格闭子概形"的一般证明简单。使用自己并不理解的证明所得的结果，有时是必要的，但终究危险，我们不想鼓励这种习惯。遗憾的是，由于时间有限，我们接受了在这一篇文章中对几个次要之点草率处理。}

设 $\gamma\in G^{\mathrm{rs}}(F^{\mathrm{sep}})$，其像为 $a\in A^{\mathrm{rs}}(F^{\mathrm{sep}})$，$T=G_\gamma$ 是其中心化子。有切空间的短正合列
\[
0\to\operatorname{Tan}_\gamma\bigl(c^{-1}(a)\bigr)\to\operatorname{Tan}_\gamma G\to\operatorname{Tan}_b(A)\to0
\]
\yizhu{原文此处写 $\operatorname{Tan}_b(A)$。此处 $a$ 已是 $A$ 中的点，切空间应记为 $\operatorname{Tan}_a(A)$；下文 (3.27) 的右端也正是 $\operatorname{Tan}_a(A)$。译文照录原文。}{依据：原书 p.34–35，与 (3.27) 对照}
\srcpage{35}
它诱导出等式
\begin{equation}\tag{3.27}
\wedge^d\mathfrak g=\wedge^n\operatorname{Tan}_a(A)\otimes\wedge^{d-n}\operatorname{Tan}_\gamma\bigl(c^{-1}(a)\bigr),
\end{equation}
其中 $d=\dim(G)$，$n=\dim(A)$。于是，$G$ 上一个不变的 $d$-形式 $\omega_G$ 与 $A$ 上一个处处非零的 $n$-形式 $\omega_A$ 的选取，诱导出纤维 $c^{-1}(a)$ 上的一个 $(d-n)$-形式
\[
\omega_a=\omega_G\otimes\omega_A^{-1},
\]
它非零且 $G$-不变。

这个纤维，更确切地说是纤维上系数在 $F_v$ 中的点集，有两种描述方式。首先它是有限个 $T(F_v)\backslash G(F_v)$ 的并，由 $\{g^{-1}\gamma'g\}$ 给出，其中 $\gamma'$ 跑遍与 $\gamma$ 相联系的稳定共轭类中各共轭类的代表；其次它是 $a$ 在 $c$ 下的原像。第一种形式下的传统测度，如公式 (3.17) 那样，由局部归一化测度的商 $d\bar g_v=dg_v/dt_v$ 给出。而在 $a$ 的原像上，无论局部还是整体，都有形式 $\omega_a$ 与相联系的测度 $|\omega_a|_v$。这两个测度并不相同。在命题 5.6 的证明中，我们宁愿使用第二个，它更几何。但传统上使用的是第一个，它在调和分析、从而在轨道积分理论中也更有用；这一理论由 Harish-Chandra 创立，Shalika 与（在内窥框架下）Shelstad 也作出了重要贡献。由于我们后面要用到这一理论，就必须理解两个测度之间的关系，而这个关系相当简单。

由于两个测度都相对于中心 $Z$ 不变，只需处理 $G=G_{\mathrm{der}}$ 的情形，按我们的假设它是单连通的。设 $r=r_{\mathrm{der}}$ 是 $G_{\mathrm{der}}$ 的秩，$\xi_1,\dots,\xi_r$ 是 3.3 节中表示 $\rho_i$ 的支配权。设 $\omega_T=d\xi_1\wedge\cdots\wedge d\xi_r$，$\omega_{T\backslash G}$ 是一个右不变补形式。它决定了测度 $d\bar g_v$。令 $\omega_G=\omega_T\wedge\omega_{T\backslash G}$。对 $\omega_A$，在我们的假设下它是 $B$ 上的形式，我们取
\begin{equation}\tag{3.28}
\omega_A=\omega_B=db_1\wedge\cdots\wedge db_r,
\end{equation}
这里 $b_i$ 不加区分地表示 $T$ 上与 $B$ 上的函数。

映射 $T=G_x\to A$ 在 $A^{\mathrm{rs}}$ 上是平展的。$\omega_A$ 的原像显然是 $db_1\wedge\cdots\wedge db_r$。为 $G_x$ 的特征固定一个序。

\setcounter{theorem}{28}
\begin{proposition}
若在根上选取一个序，并令
\[
\Delta(t)=\pm t^{-\rho}\prod_{\xi>0}\bigl(\xi(t)-1\bigr),
\]
\srcpageinl{36}
则
\[
\omega_A=db_1\wedge\cdots\wedge db_r=\pm\Delta(t)\,d\xi_1\wedge\cdots\wedge d\xi_r=\pm\Delta(t)\,\omega_T.
\]
这个等式中的符号无关紧要。$\Delta(t)$ 的平方是 $D(t)=\prod_\xi(\xi(t)-1)$，这是一个不变函数，从而是 $A$ 与 $B$ 上的函数。我们也把它写作 $D(a)$。函数 $\Delta(a)$ 只在相差一个符号的意义下确定，但我们使用的是范数 $|\Delta(a)|_v$，它定义良好。\end{proposition}

在证明命题之前，先计算 $\omega_a$。为此须固定 $\omega_G$ 与 $\omega_A$。这就以任意方式固定了局部测度，不要紧，因为没有任何整体效应。测度 $\omega_A$ 与 $\omega_T$ 是给定的。映射 $(\gamma,g)\mapsto g^{-1}\gamma g$ 从 $T\times T\backslash G$ 到 $G$，给出切空间之间的映射。按意义自明的记号，令
\[
\omega_{T\backslash G}=\wedge_\xi d\xi,\qquad
\omega_G=\omega_T\wedge\omega_{T\backslash G},
\]
其中 $\xi$ 跑遍 $G$ 的根。于是映射 $t\times\bar g\mapsto g^{-1}tg=t(t^{-1}g^{-1}tg)$ 从 $T\times T\backslash G$ 到 $G$，在 $T^{\mathrm{rs}}\times T\backslash G$ 上平展，而 $\omega_T\wedge\omega_{T\backslash G}$ 的像为
\[
\Bigl\{\prod_\xi\bigl(\xi(t)-1\bigr)\Bigr\}\omega_T\wedge\omega_{T\backslash G}
=\Delta^2(t)\,\omega_T\wedge\omega_{T\backslash G}=\Delta^2(t)\,\omega_G.
\]
于是由命题得到
\begin{equation}\tag{3.30}
\omega_a=\pm\Delta(t)\,\omega_{T\backslash G}.
\end{equation}
由这个等式推出关系
\begin{equation}\tag{3.31}
\int_{c^{-1}(a)}f_v(g)\,|\omega_a|_v
=|\Delta(t)|_v\,L_v(1,\sigma_{T/G})\,\Orb(t_{\mathrm{st}},f_v),
\end{equation}
其中 $c(t)=a$。左边的积分从几何角度看有它的优点，但传统上在局部调和分析中使用的是右边的表达式。等式 (3.31) 在按上述方式选取测度时成立，以下我们作此假设。因此左边用的是几何测度，而右边用的是局部归一化测度。人很容易在这两者之间迷失！非阿贝尔调和分析的专家习惯于右边的表达式，几何学家习惯于左边的。对导出群单连通的群 $G$，这个等式仍然成立。我们令
\begin{equation}\tag{3.32}
\theta_v(a;s)=L_v(s,\sigma_{T/G})\,|\Delta(t)|_v\,\Orb(t_{\mathrm{st}},f_v),
\qquad s\ge1,\quad c(t)=a,
\end{equation}

\srcpage{37}
并注意：若 $t=\gamma\in T(F)$ 是正则元，则表达式 (3.25) 等于 $\theta(a;s)=\prod_v\theta_v(a;s)$，其中 $\gamma\mapsto a$，这是因为 $\prod_v|\Delta(\gamma)|=1$。若 $\gamma$ 不正则，它等于 $0$。

我们指出，$\theta_v(a;s)$ 的局部行为，从而其 Fourier 变换——这一研究将在 [L5] 中开始——的整体行为，深受因子 $L_v(s,\sigma_{T/G})$ 的影响，因为这些函数的值依赖于 $T$ 的类，而 $T$ 又依赖于 $a$。

我们冒昧在复数域上给出这个命题的证明。任意域上的命题可以由此容易地推出。命题 3.29 中符号的任意性是显然的，它依赖于根上序的选取。同样显然的是，$\Delta$ 是特征 $t^\lambda$ 的整系数线性组合，因为每个 $a_i$ 都是。\yizhu{原文此处写 $a_i$，下文又说「诸 $a_i$ 是对称函数」。按上下文，此处的 $a_i$ 当为 Steinberg--Hitchin 基的坐标函数 $b_i$；译文照录原文。}{依据：原书 p.37；对照 (3.13)(3.34) 中 $b_i$ 的地位}这些整数与域 $F$ 无关。因此只需在 $\mathbb C$ 上验证公式。此时 $\Delta(t)$ 显然是 $t$ 的交错函数。若 $s$ 是 Weyl 群的一个反射，则 $\Delta(s(t))=-\Delta(t)$，因为诸 $a_i$ 是对称函数，而 $s$ 作为特征空间上的线性映射其行列式为 $-1$。由此 $\Delta(t)$ 在 $T$ 中的簇 $t^\alpha=1$ 上消没。于是
\begin{equation}\tag{3.33}
\Delta(t)=\pm Q(t)\,t^{-\rho}\prod_{\alpha>0}(1-t^\alpha),
\end{equation}
其中 $Q(t)$ 是函数 $t^\lambda$（$\lambda$ 为 $T$ 的特征）的线性组合。此外，若 $s$ 是与单根 $\beta$ 相联系的反射，则
\[
s\Bigl(\prod_{\alpha>0,\ \alpha\ne\beta}(1-t^\alpha)\Bigr)
=\prod_{\alpha>0,\ \alpha\ne\beta}(1-t^\alpha),
\]
且
\[
s(t^\beta)=t^{-\beta},\qquad s(t^{-\rho})=t^{\beta-\rho},
\]
于是 $Q(t)$ 是 Weyl 群下的不变函数。因此
\begin{equation}\tag{3.34}
Q(t)=\sum_{\lambda\ge0}a_\lambda S_\lambda(t),
\end{equation}
其中 $S_\lambda$ 就是 $t^{\lambda'}$ 在 $\lambda$ 的所有共轭 $\lambda'$ 上的和。取使 $a_\lambda\ne0$ 的最大 $\lambda$。于是 (3.33) 右端的最高权是 $\lambda+\rho$。由于每个 $b_i$ 都是类似于 (3.34) 的和，只是其中使 $a_\lambda\ne0$ 的最高权是基本权 $\mu_i$。

由于在这个证明中基域已变成复数域，我们混用了加法与乘法记号：$\xi_i=e^{\mu_i(z)}$，其中 $z=\sum_iz_i\check\alpha_i$ 属于 $T$ 的 Lie 代数。于是
\[
d\xi_1\wedge\cdots\wedge d\xi_r=\exp\bigl(\rho(z)\bigr)\,dz_1\wedge\cdots\wedge dz_r,
\]
而
\[
db_i=\exp\bigl(\mu_i(z)\bigr)dz_i+\cdots,
\]
略去的项都具有 $\exp(\mu(z))$ 的形式，其中 $\mu<\mu_i$。由此，$\Delta(t)$ 中使 $\lambda$ 最大的项 $\exp(\lambda(z))$ 满足 $\lambda=\sum_i\mu_i=\rho$。因此 $Q(t)$ 是常数，而这个常数是 $1$。

在更系统的论述中，应当给出对任意基域 $F$ 都成立的证明。在目前阶段，我们还没有这样的雄心。
